\documentclass[11pt]{article}
\usepackage[utf8]{inputenc}
\usepackage{geometry}
\geometry{a4paper, margin=1in}
\usepackage{titlesec}
\usepackage{enumitem}
\usepackage{xcolor}

\title{\textbf{Strategic Diagnostic Snapshot™}}
\author{RankPilot Intelligence}
\date{\today}

\begin{document}

\maketitle

\section*{I. Executive Summary}
\textbf{Firm:} González de Araujo Consultores \\
\textbf{Practice Area:} Data Protection \\
\textbf{Detected Practice Model:} Standard

\section*{II. Structural Positioning}
\textbf{Current Tier Assessment:} Unranked / New Entry \\
\textbf{Strategic Confidence:} 100.0\% \\
\textbf{Advantage:} Standard Market Offerings

\section*{III. Portfolio Analysis (Key Matters)}
\begin{itemize}
\item \textbf{Grupo Hermes Data Protection Framework} (Client: Grupo Hermes | Value: N/A)\\ Grupo Hermes, a Mexican diversified infrastructure and services conglomerate, Data Protection Framework is significant because it evidences the firm’s central role in converting data protection and governance requirements into an implemented set of connected operational measures, with a quantified result in the handling of ARCO requests. Rather than addressing a single document or isolated procedure, the mandate encompassed the design and implementation of a thorough data protection and governance framework for Grupo Hermes.

The firm’s work began with mapping personal data flows, establishing the factual basis for the framework and identifying the handling of personal data relevant to the engagement. The work for Grupo Hermes, developing tailored privacy notices, implementing internal policies, and establishing procedures for ARCO rights management, formed interconnected components of the data protection and governance framework. The tailored privacy notices, internal policies and ARCO rights-management procedures were implemented alongside the personal-data-flow mapping, ensuring that the framework addressed both the governance of personal data and the operational process for handling ARCO requests.

Arturo González de Araujo and Ernesto Zavala led the engagement. Their leadership covered the design and implementation of the thorough data protection and governance framework, including the mapping of personal data flows, the development of tailored privacy notices, the implementation of internal policies and the establishment of procedures for ARCO rights management. This structured approach led to 100% compliance in handling ARCO requests and avoiding regulatory sanctions, providing a concrete and quantified outcome from the framework implemented for Grupo Hermes.
\item \textbf{Mega Direct Data Protection Advisory} (Client: MEGA DIRECT, S.A. de C.V. | Value: N/A)\\ Sixteen years of advice on data protection and database-related matters for MEGA DIRECT, S.A. de C.V., a customer experience, call center, and marketing services provider, provide sustained evidence of the team’s role in protecting the legal and operational continuity of a client’s information assets. Rather than being confined to an isolated instruction, the mandate reflects a continuing advisory relationship addressing complex questions at the point where data use, database management and business operations intersect.

MEGA DIRECT, S.A. de C.V. instructed the firm in relation to data protection and database-related matters over this 16-year period. The work included representing the client in legal proceedings concerning the acquisition, processing, storage and protection of data. Those proceedings placed the full lifecycle of the client’s data handling before a legal forum: from acquisition and processing through to storage and protection. In a highly regulated environment, the representation required the team to address the legal issues connected with those distinct stages while maintaining focus on the client’s ability to continue operating.

Arturo González de Araujo, Ernesto Zavala and Damián Ortega led the representation and the continuing advice. Their involvement connects the contentious proceedings with the wider data protection and database-related mandate, demonstrating sustained partner-level attention across 16 years. The stated outcome was business continuity for MEGA DIRECT, S.A. de C.V. in the context of legal proceedings concerning its data acquisition, processing, storage and protection activities.
\item \textbf{Biocodex Data Protection Framework} (Client: Biocodex | Value: N/A)\\ Biocodex’s Data Protection Framework engagement is significant because it translated the design of a personal data protection framework into a defined internal accountability structure for Biocodex, a global pharmaceutical company focused on probiotics and specialty healthcare. Rather than being limited to high-level personal data protection advice, the mandate focused on the design and implementation of a structured personal data protection framework for Biocodex, with implementation serving as the central feature of the work.

Arturo González de Araujo, Ernesto Zavala and Damián Ortega led the matter. Their role encompassed advising on both the framework’s design and its implementation, ensuring that the personal data protection project resulted in an identifiable institutional arrangement within Biocodex rather than remaining solely a policy-level exercise. The work therefore connected the structured framework to the client’s internal controls.

The concrete institutional outcome was the establishment of a dedicated Personal Data Protection Department. That dedicated department provided the internal structure through which the newly designed and implemented personal data protection framework could be supported and through which internal controls could be strengthened. As stated in the mandate, this departmental establishment strengthened internal controls and reduced specific regulatory risk, providing a demonstrable organisational result from the personal data protection work.

Accordingly, the Biocodex Data Protection Framework matter evidences implementation-led advice that combined framework design, framework implementation, internal controls and the creation of a dedicated Personal Data Protection Department.
\item \textbf{Hotel Riazor Data Protection Enhancement} (Client: Hotel Riazor | Value: N/A)\\ Hotel Riazor, a business hotel offering lodging and event services in Mexico City: Data Protection Enhancement was centred on embedding stronger personal-data governance within the client’s internal structures, with particular emphasis on the operational role of its Personal Data Department and the adoption of internal policies for lawful data handling.

Hotel Riazor instructed the firm to advise on strengthening its personal data protection and governance framework. Rather than treating personal-data protection as a discrete policy exercise, the mandate brought together two connected elements: reinforcing Hotel Riazor’s Personal Data Department and implementing internal policies designed to ensure lawful data handling. The work therefore focused on translating the client’s data-protection and governance objectives into internal departmental reinforcement and documented policy measures addressing the handling of personal data.

Arturo González de Araujo, Ernesto Zavala and Damián Ortega led the mandate. Their role was to advise on strengthening Hotel Riazor’s personal data protection and governance framework, while supporting the client in reinforcing its Personal Data Department and implementing the internal policies required to ensure lawful data handling. The resulting outcome was the reinforcement of Hotel Riazor’s Personal Data Department together with the implementation of internal policies directed at ensuring lawful data handling.
\item \textbf{Grupo Excelsior Data Protection Regularisation} (Client: Grupo Excelsior | Value: N/A)\\ Regularisation and restructuring of operations can require data protection to be addressed as an integral operational workstream rather than as an isolated policy exercise. In this mandate, Grupo Excelsior, one of Mexico’s leading dairy producers, engaged the team to advise on the regularisation and restructuring of operations focusing on data protection, placing the design of internal governance measures alongside the practical requirements of operational change.

The engagement centred on converting the data-protection focus of the regularisation exercise into implemented internal measures. Arturo González de Araujo, Ernesto Zavala and Damián Ortega led the work, designing and implementing internal policies and training programmes. Their role therefore extended beyond identifying data-protection issues: the team developed the internal policies required for the restructured operations and implemented training programmes intended to support their internal application during the regularisation process.

The resulting policies and training programmes formed a concrete part of Grupo Excelsior’s operational regularisation and restructuring. For Grupo Excelsior, significantly strengthening compliance posture and reducing regulatory exposure was the stated outcome of the team’s design and implementation work. The matter provides submission evidence of the team’s role in linking data-protection policy design, training programmes and the regularisation and restructuring of operations.
\item \textbf{Grupo Modelquipo Data Protection Framework} (Client: Grupo Modelquipo | Value: N/A)\\ Governance-level treatment of data protection is the central significance of the Grupo Modelquipo, industrial equipment and machinery solutions provider, Data Protection Framework mandate. For Grupo Modelquipo, the work was positioned not as a standalone data protection exercise, but as an element of the client’s corporate governance strategy, giving the implementation of the framework direct institutional relevance.

Grupo Modelquipo instructed the team to advise on the implementation of a data protection framework as part of that corporate governance strategy. The firm’s role was specifically to support implementation while placing data protection within the client’s governance architecture. The mandate therefore connected the practical introduction of the data protection framework with the stated objective of making data protection a core governance element.

Arturo González de Araujo, Ernesto Zavala and Damián Ortega led the work for Grupo Modelquipo. Their involvement centred on advising on the implementation of the data protection framework and its incorporation into the client’s corporate governance strategy. The reported outcome was improved regulatory compliance and the integration of data protection as a core governance element, reflecting a concrete governance-based approach to the framework’s implementation.
\item \textbf{Tiendas Chedraui Data Protection Advisory} (Client: Tiendas Chedraui | Value: N/A)\\ Tiendas Chedraui's significance in embedding data-protection risk assessment into the legal review of business growth, so that new store openings and operational expansion are assessed with regulatory requirements and governance standards in view from the outset.

Tiendas Chedraui, a Mexican supermarket and retail chain, has instructed the firm in connection with the legal assessment of new store openings and operational expansion. The engagement is directed specifically at the data-processing implications arising as the client opens new stores and develops its operations. Rather than treating data protection as a separate or abstract compliance exercise, the mandate connects the evaluation of data-processing risks directly to the client’s proposed expansion activities and the legal assessment required for each development.

Arturo González de Araujo and Damián Ortega lead the work. Their role is to evaluate the risks associated with data processing in the context of Tiendas Chedraui’s new store openings and operational expansion, and to ensure alignment with regulatory requirements and governance standards. The resulting focus is a concrete legal assessment of how data-processing risks should be considered as the client advances new openings and expands its operations.

\end{itemize}

\section*{IV. Strategic Evolution Path}
\begin{itemize}
\item Maintain current strategy.
\end{itemize}

\vfill
\centering
\small{This document is a structural assessment and does not guarantee ranking results.}

\end{document}